
\documentclass[11pt]{article}
\usepackage[a4paper,margin=1in]{geometry}
\usepackage{amsmath,amssymb,amsfonts,bm}
\usepackage{mathtools}
\usepackage{hyperref}

\title{Monolithic one--domain biofilm detachment \& growth:\\
Diffuse--interface Navier--Stokes--Brinkman--Biot formulation\\
with primal variables, weak form, residual, and consistent Jacobian}
\author{}
\date{}

% ---------- macros ----------
\newcommand{\R}{\mathbb{R}}
\newcommand{\dd}{\,\mathrm{d}}
\newcommand{\I}{\bm{I}}
\newcommand{\n}{\bm{n}}
\newcommand{\grad}{\nabla}
\newcommand{\divv}{\nabla\!\cdot}
\newcommand{\eps}{\bm{\varepsilon}}
\newcommand{\sym}{\operatorname{sym}}
\newcommand{\tr}{\operatorname{tr}}
\newcommand{\inner}[2]{\left(#1,#2\right)}
\newcommand{\ip}[2]{\left\langle #1,#2\right\rangle}
\newcommand{\Dt}{\partial_t}
\newcommand{\jump}[1]{\llbracket #1\rrbracket}
\newcommand{\avg}[1]{\{#1\}}
\newcommand{\Hs}{H_\varepsilon}
\newcommand{\ds}{\delta_\varepsilon}

\begin{document}
\maketitle

\section{Goal and modeling choice}

We want a formulation that
(i) reduces to incompressible Navier--Stokes in the free fluid,
(ii) reduces to a poroelastic biofilm (Biot--type) in the biofilm,
(iii) allows \emph{growth and detachment} without remeshing, by letting a smooth indicator field move the material distribution,
and (iv) is suitable for a \emph{monolithic Newton solve} (residual + consistent Jacobian).

The essential ingredient is a \emph{biofilm indicator} (diffuse interface) field
\[
\alpha^{bf}(x,t)\in[0,1],
\]
where we adopt the convention
\[
\alpha^{bf}=1\ \text{in biofilm},\qquad \alpha^{bf}=0\ \text{in pure fluid},
\]
and $\alpha^{bf}$ varies smoothly across a thin layer of thickness $\varepsilon$.

\paragraph{Two different ``fractions'' appear and should not be conflated.}
\begin{itemize}
\item $\alpha^{bf}$: \emph{where biofilm exists} (domain indicator / phase field).
\item $\phi$: \emph{porosity inside the biofilm} (fluid volume fraction of the \emph{biofilm material}).
In pure fluid we enforce $\phi\approx 1$.
\end{itemize}

\section{Smooth partition and diffuse interface operators}

Define smooth characteristic weights for the two regions
\[
\chi_B(\alpha)=\alpha,\qquad \chi_F(\alpha)=1-\alpha.
\]
(You may replace these by a smoothed Heaviside $\Hs(\alpha)$ if you use a double-well phase field; the algebra is identical.)

If you still want ``interface'' terms (e.g. a Nitsche-like stabilization),
you can represent surface integrals over the moving interface $\Gamma(t)$ by diffuse-interface identities
\[
\int_{\Gamma(t)} g\,\dd s \;\approx\; \int_{\Omega} g \,\ds(\alpha)\,\dd x,
\qquad
\ds(\alpha)\approx |\grad \alpha| \quad \text{(or } \ds(\alpha)=4\alpha(1-\alpha)/\varepsilon\text{)}.
\]
A diffuse normal can be defined by
\[
\n(\alpha)=\frac{\grad\alpha}{|\grad\alpha|+\eta},\qquad \eta\ll 1.
\]

\section{Primal variables and kinematics}

Let $\Omega\subset\R^d$ be a fixed Eulerian computational domain, $d=2$ or $3$.

\subsection*{Primal (unknown) fields}

We solve \emph{one} monolithic system for the vector of unknowns
\[
\boxed{
\bm{w} = \big(\bm{v},\,p,\,\bm{u},\,\phi,\,\alpha,\,S\big),
}
\]
where
\begin{itemize}
\item $\bm{v}(x,t)$: fluid velocity (free-fluid velocity when $\alpha\approx 0$, pore-fluid velocity when $\alpha\approx 1$);
\item $p(x,t)$: pressure (free-fluid pressure when $\alpha\approx 0$, pore pressure when $\alpha\approx 1$);
\item $\bm{u}(x,t)$: skeleton displacement (active only where $\alpha\approx 1$);
\item $\phi(x,t)\in(0,1]$: porosity of the biofilm material (enforced $\phi\to 1$ when $\alpha\to 0$);
\item $\alpha(x,t)\in[0,1]$: biofilm indicator (growth/detachment moves $\alpha$);
\item $S(x,t)\ge 0$: substrate/nutrient concentration for Monod growth.
\end{itemize}

\subsection*{Skeleton velocity and acceleration}
We denote the skeleton velocity by
\[
\bm{v}^S := \Dt \bm{u},
\]
and (if needed) its material acceleration in Eulerian form by
\[
\bm{a}^S := \Dt \bm{v}^S + (\bm{v}^S\cdot\grad)\bm{v}^S.
\]
In many biofilm problems one sets the skeleton to quasi-static ($\bm{a}^S=0$); the derivations below cover both cases.

\subsection*{Small-strain vs large-deformation skeleton}
For compactness we write the skeleton stress in Cauchy form as $\bm{\sigma}_{eff}(\bm{u})$.
You may choose:
\begin{itemize}
\item small strain: $\bm{\sigma}_{eff}=2\mu_s \eps(\bm{u})+\lambda_s \tr(\eps(\bm{u}))\I$, $\eps(\bm{u})=\sym(\grad \bm{u})$;
\item large deformation (Eulerian reference-map): $F=(I-\grad u)^{-1}$ etc.\ (as in your current fully Eulerian poroelastic module) and $\bm{\sigma}_{eff}$ neo-Hookean.
\end{itemize}
Only the Gateaux derivative $D\bm{\sigma}_{eff}[\delta\bm{u}]$ is needed for the Jacobian.

\section{Constitutive interpolation (turn physics on/off with $\alpha$)}

We define \emph{effective} material coefficients via $\alpha$:
\begin{align*}
\rho(\alpha,\phi) &:= \chi_F(\alpha)\,\rho_f + \chi_B(\alpha)\,\rho_f \phi, \\
\mu(\alpha,\phi) &:= \chi_F(\alpha)\,\mu_f + \chi_B(\alpha)\,\mu_B(\phi),
\end{align*}
where $\mu_B(\phi)$ is the Brinkman viscosity in the porous region (often $\mu_B=\mu_f$ or $\mu_B=\phi\mu_f$).

The Darcy/Brinkman drag coefficient is activated in biofilm:
\[
\boxed{
\beta(\alpha,\phi,\bm{u}) := \chi_B(\alpha)\,\mu_f\,\phi^2\,\kappa(\phi,\bm{u})^{-1},
}
\]
with permeability $\kappa$ (possibly deformation-dependent). If you use the Eulerian reference-map kinematics, you may push forward permeability as in your existing code base.

\section{Governing equations (strong form)}

\subsection{Generalized Navier--Stokes--Brinkman momentum (whole domain)}
\begin{equation}\label{eq:momentum}
\rho(\alpha,\phi)\Big(\Dt \bm{v} + (\bm{v}\cdot\grad)\bm{v}\Big)
-\divv\big(2\mu(\alpha,\phi)\,\eps(\bm{v})\big) + \grad p
+ \beta(\alpha,\phi,\bm{u})\big(\bm{v}-\bm{v}^S\big)
= \bm{f}_v
\quad \text{in }\Omega.
\end{equation}
When $\alpha\to 0$, $\beta\to 0$ and $\mu\to\mu_f$, so \eqref{eq:momentum} becomes Navier--Stokes.
When $\alpha\to 1$ and $\beta$ is large (small permeability), \eqref{eq:momentum} reduces to Brinkman/Darcy flow relative to the moving skeleton.

\subsection{Saturation / volume constraint (blends fluid and porous incompressibility)}
We enforce divergence-free in pure fluid and the correct mixture constraint in biofilm:
\begin{equation}\label{eq:mass}
\divv\Big(\chi_F(\alpha)\,\bm{v} + \chi_B(\alpha)\big(\phi\bm{v}+(1-\phi)\bm{v}^S\big)\Big)
= \chi_B(\alpha)\,s_v
\quad \text{in }\Omega.
\end{equation}
Here $s_v$ represents net volumetric source from growth/decay if you do not explicitly model the mass of consumed substrate.

\subsection{Skeleton momentum (poroelasticity, active where $\alpha\approx 1$)}
A Biot-type skeleton balance with effective stress $\bm{\sigma}_{eff}(\bm{u})$ is:
\begin{equation}\label{eq:skeleton}
\chi_B(\alpha)\Big(
\rho_s(1-\phi)\,\bm{a}^S
-\divv \bm{\sigma}_{eff}(\bm{u})
+(1-\phi)\grad p
-\beta(\alpha,\phi,\bm{u})\big(\bm{v}-\bm{v}^S\big)
\Big)
= \chi_B(\alpha)\,\bm{f}_u
\quad \text{in }\Omega.
\end{equation}
Outside the biofilm ($\alpha\approx 0$) the equation vanishes (no skeleton).

\subsection{Porosity evolution (biofilm material law, active where $\alpha\approx 1$)}
Assuming incompressible solid constituent with true density $\rho_s^\ast$,
a consistent Eulerian porosity evolution is
\begin{equation}\label{eq:phi_evol}
\chi_B(\alpha)\Big(
\Dt\phi + \bm{v}^S\cdot\grad\phi - (1-\phi)\divv \bm{v}^S
+ \frac{\Pi_b(S,\phi,\alpha,\tau)}{\rho_s^\ast}
\Big)=0
\quad \text{in }\Omega,
\end{equation}
where $\Pi_b$ is the (signed) biofilm mass source per mixture volume:
positive for growth, negative for decay/erosion.
A Monod choice is
\[
\Pi_b = \rho_s^\ast\,\mu_{\max}\frac{S}{K_S+S}\,(1-\phi)\,\alpha - \rho_s^\ast k_d (1-\phi)\,\alpha.
\]
(You can also include detachment as a sink in $\Pi_b$; alternatively treat detachment through the $\alpha$-equation below.)

To enforce $\phi\to 1$ in pure fluid, add a relaxation:
\[
(1-\alpha)\,(\phi-1)=0
\quad\text{(weakly, via a penalty in the weak form)}.
\]

\subsection{Biofilm indicator evolution (growth + detachment + diffuse interface)}
A practical monolithic evolution for $\alpha$ is an advection--diffusion--reaction equation:
\begin{equation}\label{eq:alpha_evol}
\Dt \alpha + \bm{v}^S\cdot\grad\alpha
= G(S,\phi)\,\alpha(1-\alpha) \;-\; D(\tau)\,\alpha
\;+\; \divv(D_\alpha \grad\alpha)
\quad \text{in }\Omega.
\end{equation}
Here
\[
G(S,\phi)=k_g\,\mu_{\max}\frac{S}{K_S+S}\,(1-\phi),
\qquad
D(\tau)=k_{\det}\,g(\tau),
\]
with $\tau$ a shear measure computed from the fluid stress, e.g.
\[
\tau = \|\bm{P}_\tau\big(2\mu_f\eps(\bm{v})\n\big)\|,\qquad \bm{P}_\tau = \I - \n\otimes\n.
\]
The logistic factor $\alpha(1-\alpha)$ keeps $\alpha$ in $[0,1]$; if you want sharper interfaces, add a double-well term $-M W'(\alpha)$ (Allen--Cahn).

\subsection{Substrate transport (whole domain, with correct biofilm capacity)}
A conservative choice that treats substrate in free fluid and in biofilm pore space is:
\begin{equation}\label{eq:substrate}
\Dt\big(\chi_F(\alpha)S + \chi_B(\alpha)\phi S\big)
+\divv\big(\chi_F(\alpha)S\bm{v} + \chi_B(\alpha)\phi S \bm{v}\big)
= \divv\big(D_S(\alpha,\phi)\grad S\big) - R_S,
\end{equation}
with consumption tied to growth, e.g.\ $R_S = \frac{1}{Y}\Pi_b/\rho_s^\ast$ (plus maintenance terms).

\section{Weak form (monolithic)}

Let the trial spaces be
\[
\bm{v}\in \bm{V},\ p\in Q,\ \bm{u}\in \bm{U},\ \phi\in W_\phi,\ \alpha\in W_\alpha,\ S\in W_S,
\]
and test functions
\[
\bm{w}\in \bm{V}_0,\ q\in Q,\ \bm{\eta}\in \bm{U}_0,\ \zeta\in W_\phi,\ \xi\in W_\alpha,\ r\in W_S.
\]
(Choose boundary conditions as usual; the weak form below only shows interior contributions.)

Denote backward Euler time discretization:
\[
\Dt \bm{v}\approx \frac{\bm{v}-\bm{v}^n}{\Delta t},\quad
\Dt \bm{u}\approx \frac{\bm{u}-\bm{u}^n}{\Delta t},\quad
\Dt \phi\approx \frac{\phi-\phi^n}{\Delta t},\quad
\Dt \alpha\approx \frac{\alpha-\alpha^n}{\Delta t},\quad
\Dt S\approx \frac{S-S^n}{\Delta t}.
\]
Let $\bm{v}^S=\frac{\bm{u}-\bm{u}^n}{\Delta t}$.

\subsection*{Residual functional}
Define the residual $R(\bm{v},p,\bm{u},\phi,\alpha,S;\bm{w},q,\bm{\eta},\zeta,\xi,r)$ as the sum:
\[
R = R_v + R_p + R_u + R_\phi + R_\alpha + R_S.
\]

\paragraph{(i) Momentum residual $R_v$.}
\begin{align}
R_v :=
\int_\Omega &
\rho(\alpha,\phi)\Big(\frac{\bm{v}-\bm{v}^n}{\Delta t}+(\bm{v}\cdot\grad)\bm{v}\Big)\cdot \bm{w}
+ 2\mu(\alpha,\phi)\,\eps(\bm{v}):\eps(\bm{w})
- p\,\divv \bm{w} \notag\\
&\quad
+ \beta(\alpha,\phi,\bm{u})\big(\bm{v}-\bm{v}^S\big)\cdot \bm{w}
- \bm{f}_v\cdot\bm{w}
\ \dd x.
\label{eq:Rv}
\end{align}

\paragraph{(ii) Mass/volume residual $R_p$.}
\begin{equation}
R_p :=
\int_\Omega
q\Big(
\divv\big(\chi_F(\alpha)\bm{v}+\chi_B(\alpha)(\phi\bm{v}+(1-\phi)\bm{v}^S)\big)
-\chi_B(\alpha)s_v
\Big)\dd x.
\label{eq:Rp}
\end{equation}

\paragraph{(iii) Skeleton residual $R_u$.}
\begin{align}
R_u :=
\int_\Omega &
\chi_B(\alpha)\Big[
\rho_s(1-\phi)\,\bm{a}^S\cdot\bm{\eta}
+ \bm{\sigma}_{eff}(\bm{u}):\grad\bm{\eta}
-(1-\phi)p\,\divv\bm{\eta}
-\beta(\alpha,\phi,\bm{u})\big(\bm{v}-\bm{v}^S\big)\cdot\bm{\eta}
-\bm{f}_u\cdot\bm{\eta}
\Big]\dd x.
\label{eq:Ru}
\end{align}
If quasi-static, set $\bm{a}^S=\bm{0}$.

\paragraph{(iv) Porosity residual $R_\phi$.}
Add diffusion/regularization of $\phi$ (optional but practical), and enforce $\phi\to 1$ when $\alpha\to 0$:
\begin{align}
R_\phi := \int_\Omega &
\chi_B(\alpha)\,\zeta\Big(
\frac{\phi-\phi^n}{\Delta t} + \bm{v}^S\cdot\grad\phi - (1-\phi)\divv\bm{v}^S
+ \frac{\Pi_b(S,\phi,\alpha,\tau)}{\rho_s^\ast}
\Big)\dd x \notag\\
&+ \int_\Omega D_\phi\,\grad\phi\cdot\grad\zeta\ \dd x
+ \int_\Omega \gamma_\phi(1-\alpha)\,(\phi-1)\,\zeta\ \dd x.
\label{eq:Rphi}
\end{align}

\paragraph{(v) Biofilm indicator residual $R_\alpha$.}
\begin{align}
R_\alpha := \int_\Omega &
\xi\Big(
\frac{\alpha-\alpha^n}{\Delta t} + \bm{v}^S\cdot\grad\alpha
- G(S,\phi)\alpha(1-\alpha) + D(\tau)\alpha
\Big)\dd x
+ \int_\Omega D_\alpha\,\grad\alpha\cdot\grad\xi\ \dd x.
\label{eq:Ralpha}
\end{align}

\paragraph{(vi) Substrate residual $R_S$.}
Let $C(\alpha,\phi):=\chi_F(\alpha)+\chi_B(\alpha)\phi$ be the local capacity.
\begin{align}
R_S := \int_\Omega &
r\Big(
\frac{C(\alpha,\phi)S - C(\alpha^n,\phi^n)S^n}{\Delta t}
+ \divv\big(C(\alpha,\phi)S\bm{v}\big)
- \divv(D_S(\alpha,\phi)\grad S)
+ R_S(S,\phi,\alpha)
\Big)\dd x.
\label{eq:RS}
\end{align}

\section{Consistent Newton linearization: Jacobian}

Let the Newton increment be
\[
\delta\bm{w} = (\delta\bm{v},\delta p,\delta\bm{u},\delta\phi,\delta\alpha,\delta S).
\]
The consistent Jacobian is the Gateaux derivative
\[
\boxed{
J(\delta\bm{w};\bm{w},q,\bm{\eta},\zeta,\xi,r)
:= D R(\bm{v},p,\bm{u},\phi,\alpha,S)\,[\delta\bm{w}],
}
\]
and the Newton step solves $J(\delta\bm{w}) = -R$ for all tests.

Below, we list the Jacobian contributions block-by-block. For brevity we denote $\Delta t$ by $dt$ and $\bm{v}^S=\frac{\bm{u}-\bm{u}^n}{dt}$, so $\delta \bm{v}^S=\frac{\delta\bm{u}}{dt}$.

\subsection{Jacobian of momentum residual $R_v$}

\paragraph{Derivative w.r.t.\ $\bm{v}$.}
\begin{align}
D_{\bm{v}}R_v[\delta\bm{v}] &=
\int_\Omega
\rho(\alpha,\phi)\Big(\frac{\delta\bm{v}}{dt} + (\delta\bm{v}\cdot\grad)\bm{v} + (\bm{v}\cdot\grad)\delta\bm{v}\Big)\cdot\bm{w}
+ 2\mu(\alpha,\phi)\,\eps(\delta\bm{v}):\eps(\bm{w})
\notag\\
&\qquad
+ \beta(\alpha,\phi,\bm{u})\,\delta\bm{v}\cdot\bm{w}\ \dd x.
\end{align}

\paragraph{Derivative w.r.t.\ $p$.}
\[
D_{p}R_v[\delta p] = \int_\Omega -\delta p\,\divv\bm{w}\ \dd x.
\]

\paragraph{Derivative w.r.t.\ $\bm{u}$ (through $\bm{v}^S$ and possibly $\beta$).}
\begin{align}
D_{\bm{u}}R_v[\delta\bm{u}]
&=\int_\Omega
\beta(\alpha,\phi,\bm{u})\Big(-\delta\bm{v}^S\Big)\cdot\bm{w}\ \dd x
\;+\;\int_\Omega \Big(D_{\bm{u}}\beta[\delta\bm{u}]\Big)\big(\bm{v}-\bm{v}^S\big)\cdot\bm{w}\ \dd x,
\end{align}
with $\delta\bm{v}^S=\delta\bm{u}/dt$.

\paragraph{Derivative w.r.t.\ $\phi$ and $\alpha$ (variable coefficients).}
Using $\rho_\phi=\partial\rho/\partial\phi$, $\rho_\alpha=\partial\rho/\partial\alpha$, similarly for $\mu$ and $\beta$:
\begin{align}
D_{\phi}R_v[\delta\phi] &=
\int_\Omega
\rho_\phi\,\delta\phi\Big(\frac{\bm{v}-\bm{v}^n}{dt}+(\bm{v}\cdot\grad)\bm{v}\Big)\cdot\bm{w}
+2\mu_\phi\,\delta\phi\,\eps(\bm{v}):\eps(\bm{w})
+\beta_\phi\,\delta\phi\,(\bm{v}-\bm{v}^S)\cdot\bm{w}\ \dd x,
\\
D_{\alpha}R_v[\delta\alpha] &=
\int_\Omega
\rho_\alpha\,\delta\alpha\Big(\frac{\bm{v}-\bm{v}^n}{dt}+(\bm{v}\cdot\grad)\bm{v}\Big)\cdot\bm{w}
+2\mu_\alpha\,\delta\alpha\,\eps(\bm{v}):\eps(\bm{w})
+\beta_\alpha\,\delta\alpha\,(\bm{v}-\bm{v}^S)\cdot\bm{w}\ \dd x.
\end{align}

\subsection{Jacobian of mass/volume residual $R_p$}

Let $F_m(\bm{v},\bm{u},\phi,\alpha)=\chi_F(\alpha)\bm{v}+\chi_B(\alpha)(\phi\bm{v}+(1-\phi)\bm{v}^S)$.
Then $R_p=\int q(\divv F_m - \chi_B s_v)\dd x$.

\paragraph{Derivative w.r.t.\ $\bm{v}$ and $\bm{u}$.}
\begin{align}
D_{\bm{v}}R_p[\delta\bm{v}] &= \int_\Omega q\,\divv\Big(\chi_F(\alpha)\delta\bm{v}+\chi_B(\alpha)\phi\,\delta\bm{v}\Big)\dd x,\\
D_{\bm{u}}R_p[\delta\bm{u}] &= \int_\Omega q\,\divv\Big(\chi_B(\alpha)(1-\phi)\delta\bm{v}^S\Big)\dd x, \qquad \delta\bm{v}^S=\delta\bm{u}/dt.
\end{align}

\paragraph{Derivative w.r.t.\ $\phi$ and $\alpha$.}
\begin{align}
D_{\phi}R_p[\delta\phi] &= \int_\Omega q\,\divv\Big(\chi_B(\alpha)\,\delta\phi\,(\bm{v}-\bm{v}^S)\Big)\dd x,\\
D_{\alpha}R_p[\delta\alpha] &= \int_\Omega q\,\divv\Big(\chi_F'(\alpha)\delta\alpha\,\bm{v}+\chi_B'(\alpha)\delta\alpha(\phi\bm{v}+(1-\phi)\bm{v}^S)\Big)\dd x
-\int_\Omega q\,\chi_B'(\alpha)\delta\alpha\,s_v\ \dd x.
\end{align}
With $\chi_F(\alpha)=1-\alpha$, $\chi_B(\alpha)=\alpha$, we have $\chi_F'=-1$, $\chi_B'=1$.

\subsection{Jacobian of skeleton residual $R_u$}

\paragraph{Derivative w.r.t.\ $\bm{u}$.}
This includes inertia (if used), effective stress tangent, and coupling through $\bm{v}^S$:
\begin{align}
D_{\bm{u}}R_u[\delta\bm{u}] &=
\int_\Omega \chi_B(\alpha)\Big[
\rho_s(1-\phi)\,D\bm{a}^S[\delta\bm{u}]\cdot\bm{\eta}
+ \big(D\bm{\sigma}_{eff}[\delta\bm{u}]\big):\grad\bm{\eta}
-(1-\phi)p\,\divv(\delta\bm{\eta})\cdot 0
+\beta\,\delta\bm{v}^S\cdot\bm{\eta}
\Big]\dd x
\\
&\qquad
-\int_\Omega \chi_B(\alpha)\Big(D_{\bm{u}}\beta[\delta\bm{u}]\Big)(\bm{v}-\bm{v}^S)\cdot\bm{\eta}\ \dd x,
\end{align}
where $\delta\bm{v}^S=\delta\bm{u}/dt$ and $D\bm{a}^S$ is the chosen acceleration linearization.
If quasi-static, drop the $\bm{a}^S$ terms.

\paragraph{Derivative w.r.t.\ $\bm{v}$ and $p$.}
\[
D_{\bm{v}}R_u[\delta\bm{v}] = \int_\Omega \chi_B(\alpha)\Big(-\beta\,\delta\bm{v}\cdot\bm{\eta}\Big)\dd x,
\qquad
D_p R_u[\delta p] = \int_\Omega \chi_B(\alpha)\Big(-(1-\phi)\delta p\,\divv\bm{\eta}\Big)\dd x.
\]

\paragraph{Derivative w.r.t.\ $\phi$ and $\alpha$.}
\begin{align}
D_\phi R_u[\delta\phi] &=
\int_\Omega \chi_B(\alpha)\Big[
-\rho_s\,\delta\phi\,\bm{a}^S\cdot\bm{\eta}
+\delta\phi\,p\,\divv\bm{\eta}
-\beta_\phi\,\delta\phi\,(\bm{v}-\bm{v}^S)\cdot\bm{\eta}
\Big]\dd x,
\\
D_\alpha R_u[\delta\alpha] &=
\int_\Omega \chi_B'(\alpha)\delta\alpha\Big[
\rho_s(1-\phi)\bm{a}^S\cdot\bm{\eta}
+\bm{\sigma}_{eff}(\bm{u}):\grad\bm{\eta}
-(1-\phi)p\,\divv\bm{\eta}
-\beta(\bm{v}-\bm{v}^S)\cdot\bm{\eta}
-\bm{f}_u\cdot\bm{\eta}
\Big]\dd x.
\end{align}

\subsection{Jacobian of porosity residual $R_\phi$}

Let
\[
F_\phi := \frac{\phi-\phi^n}{dt} + \bm{v}^S\cdot\grad\phi - (1-\phi)\divv\bm{v}^S + \frac{\Pi_b}{\rho_s^\ast}.
\]
Then $R_\phi = \int \chi_B(\alpha)\zeta F_\phi + D_\phi \grad\phi\cdot\grad\zeta + \gamma_\phi(1-\alpha)(\phi-1)\zeta$.

Key derivatives:
\begin{align}
D_\phi R_\phi[\delta\phi] &=
\int_\Omega \chi_B(\alpha)\zeta\Big(\frac{\delta\phi}{dt} + \bm{v}^S\cdot\grad\delta\phi + (\delta\phi)\divv\bm{v}^S + \frac{\partial_\phi \Pi_b}{\rho_s^\ast}\delta\phi\Big)\dd x
+ \int_\Omega D_\phi \grad\delta\phi\cdot\grad\zeta\ \dd x
+ \int_\Omega \gamma_\phi(1-\alpha)\delta\phi\,\zeta\ \dd x,\\
D_{\bm{u}} R_\phi[\delta\bm{u}] &=
\int_\Omega \chi_B(\alpha)\zeta\Big( \delta\bm{v}^S\cdot\grad\phi - (1-\phi)\divv\delta\bm{v}^S + \frac{\partial_{\bm{u}}\Pi_b}{\rho_s^\ast}\Big)\dd x,
\quad \delta\bm{v}^S=\delta\bm{u}/dt,\\
D_\alpha R_\phi[\delta\alpha] &=
\int_\Omega \chi_B'(\alpha)\delta\alpha\,\zeta\,F_\phi\ \dd x
-\int_\Omega \gamma_\phi\,\delta\alpha\,(\phi-1)\,\zeta\ \dd x,\\
D_S R_\phi[\delta S] &= \int_\Omega \chi_B(\alpha)\zeta\,\frac{\partial_S\Pi_b}{\rho_s^\ast}\,\delta S\ \dd x.
\end{align}

\subsection{Jacobian of indicator residual $R_\alpha$}

Let
\[
F_\alpha:=\frac{\alpha-\alpha^n}{dt}+\bm{v}^S\cdot\grad\alpha - G(S,\phi)\alpha(1-\alpha) + D(\tau)\alpha.
\]
Then
\begin{align}
D_\alpha R_\alpha[\delta\alpha] &=
\int_\Omega \xi\Big(\frac{\delta\alpha}{dt}+\bm{v}^S\cdot\grad\delta\alpha
- G(S,\phi)\,(1-2\alpha)\delta\alpha + D(\tau)\delta\alpha \Big)\dd x
+\int_\Omega D_\alpha\,\grad\delta\alpha\cdot\grad\xi\ \dd x,\\
D_{\bm{u}} R_\alpha[\delta\bm{u}] &= \int_\Omega \xi\ \delta\bm{v}^S\cdot\grad\alpha\ \dd x,\\
D_\phi R_\alpha[\delta\phi] &= \int_\Omega \xi\ \Big(-\partial_\phi G(S,\phi)\,\alpha(1-\alpha)\Big)\delta\phi\ \dd x,\\
D_S R_\alpha[\delta S] &= \int_\Omega \xi\ \Big(-\partial_S G(S,\phi)\,\alpha(1-\alpha)\Big)\delta S\ \dd x.
\end{align}
If $D(\tau)$ depends on $\bm{v}$ through $\tau$, add the extra term
$D_{\bm{v}}R_\alpha[\delta\bm{v}]=\int \xi\,\alpha\,\partial_\tau D(\tau)\,D_{\bm{v}}\tau[\delta\bm{v}]\,\dd x$.

\subsection{Jacobian of substrate residual $R_S$}
This is standard. It includes derivatives of the capacity $C(\alpha,\phi)$ and any nonlinear consumption $R_S$.

\section{Optional: recovering sharp-interface Nitsche as $\varepsilon\to 0$}

If you still want Nitsche-like interface constraints (normal flux continuity, BJ slip),
you can add \emph{diffuse Nitsche} terms using $\ds(\alpha)$ and $\n(\alpha)$, e.g.
\[
R_{\Gamma,\text{pen}} = \int_\Omega \ds(\alpha)\,\frac{\gamma}{h}\,\big(\bm{v}-\bm{v}^S\big)\cdot\bm{w}\,\dd x,
\]
and similarly for tangential BJ slip with $\bm{P}_\tau(\alpha)=\I-\n(\alpha)\otimes\n(\alpha)$.
As $\varepsilon\to 0$, $\ds(\alpha)$ concentrates on $\Gamma(t)$ and these reduce to surface penalties.

\section{MMS sketch that shows detachment in time (for verification)}

Choose a moving interface encoded by $\alpha$:
\[
\alpha(x,y,t)=\frac{1}{2}\Big(1-\tanh\big(\tfrac{y-h(t)}{\varepsilon}\big)\Big),
\qquad h(t)=h_0 - V_{\det}t,
\]
so the biofilm indicator front moves downward (detachment) with speed $V_{\det}>0$.
Pick analytic fields $\bm{v},p,\bm{u},\phi,S$ (e.g.\ trigonometric-in-space, exponential-in-time),
plug them into \eqref{eq:Rv}--\eqref{eq:RS}, and \emph{define} manufactured sources $\bm{f}_v,\bm{f}_u,s_v$ and $R_S$ so that the residual is exactly zero.
Then convergence of the monolithic solver verifies your implementation including moving/diffuse detachment.

\end{document}